%==============================================================================
% Paper C -- skeleton (bullets, to be polished into prose)
% REVISED after peer review (Reviewers A & B) and an in-house leakage ablation.
%
% What changed and why (recorded for the team):
%   * The first draft led with "100% phasic-vs-steady accuracy" and explained it
%     by "phasic SNR 4.7-6.4 dB vs steady <3 dB". BOTH were wrong:
%       - steady cepstral SNR is actually 2.6-8.0 dB (4 of 6 steady workloads
%         exceed the 4.5 dB phasic gate; hashtable tops the set at 8.0). SNR does
%         NOT separate the families.
%       - an ablation (leakage_ablation.py / .json) shows the 100% is fully
%         reproduced by coverage_ratio ALONE -- a feature assigned per family
%         (0.133 steady vs 0.2 phasic). Remove the three family-conditional
%         features and accuracy falls to 0.644 (majority baseline 0.542).
%   * The paper is therefore recast from a "security detector" to a
%     "workload characterization + a cautionary leakage finding" paper.
%
% Target venue: ML-for-systems / measurement (NOT a security-detection venue
%   on this evidence).
% Build: pdflatex skeleton.tex  (IEEEtran; swap to article if missing).
%==============================================================================
\documentclass[conference]{IEEEtran}
\usepackage{amsmath,booktabs,siunitx,url}

\title{How Much Workload Identity Survives in a Single-Scalar\\
Memory-Liveness Trajectory? A Leakage-Controlled Study}

\author{\IEEEauthorblockN{TODO Author list}
\IEEEauthorblockA{TODO Affiliation}}

\begin{document}
\maketitle

\begin{abstract}
% TODO polish. Bullet stub:
\begin{itemize}
  \item Question: after collapsing per-page memory change to one scalar (APF),
        how much workload identity remains?
  \item A naive classifier reports perfect (100\%) phasic-vs-steady accuracy,
        but we show this is an artifact: a single feature that the analyzer
        computes \emph{per family} reproduces it exactly.
  \item Under a leakage-controlled evaluation -- removing every
        family-conditional feature -- the residual family signal is weak and not
        shown to beat chance: 0.644 accuracy against a 0.542 majority baseline,
        a gap that is not significant once the four non-independent replicates
        per workload-duration are accounted for.
  \item The contribution is therefore a cautionary, quantified result about
        feature-construction leakage in memory-signal classification, plus an
        honest measurement of the residual signal.
  \item On a DOF-relieved re-capture the same protocol yields a
        significant-but-coarse signal (leave-one-workload-out 0.636 vs.\ a 0.545
        baseline), showing the weak v3 result was partly a capture-fidelity
        artifact; yet the signal still fails to generalize across families,
        confirming the boundary is real, not an artifact.
\end{itemize}
\end{abstract}

\section{Introduction}
\begin{itemize}
  \item APF discards per-page detail (Paper A). A fair question is whether the
        surviving scalar still carries workload identity, at the family level
        (bursty ``phasic'' vs.\ continuous ``steady'') or the instance level
        (which of 11 workloads).
  \item \textbf{A trap we fell into, then diagnosed.} Ten APF-derived features
        plus a RandomForest separate the families perfectly under
        leave-one-replicate-out cross-validation. That number is seductive and
        \emph{wrong as a behavioral claim}: three of the ten features are
        computed conditionally on the family the analyzer already assigned.
  \item \textbf{Why this matters beyond this paper.} Memory-signal pipelines
        routinely derive ``defining metrics'' per workload class; if those
        metrics then feed a classifier, the classifier recovers the labelling
        rule, not the phenomenon. We make this concrete and measurable.
  \item \textbf{Contributions}:
    \begin{itemize}
      \item C1: a leakage-controlled evaluation protocol -- identify
            family-conditional features, ablate them, and report against a
            majority baseline with an exact confidence bound.
      \item C2: a quantified leakage demonstration -- the headline 100\% is
            reproduced by a single family-assigned feature
            (\texttt{coverage\_ratio}); the leakage-free accuracy is 0.644.
      \item C3: an honest measurement of residual identity -- a coarse scalar
            carries a \emph{weak} family signal and essentially no instance
            signal -- and the experiments needed before any security claim.
      \item C4: the residual family signal is \emph{capture-fidelity dependent}.
            Re-running the identical leakage-controlled protocol on a DOF-relieved
            capture flips the leakage-free signal from below the majority baseline
            to above it and significant -- the earlier ``not significant'' was
            partly a capture-quality artifact, not an intrinsic ceiling.
      \item C5: a generalization boundary and a regime split. The leakage-free
            signal does not transfer across \emph{families} (a held-out family's
            threats are missed; a novel subtype's family is recovered at chance),
            yet a one-class detector trained only on benign flags a novel threat
            family -- so supervised family-identification and anomaly detection
            occupy different regimes.
    \end{itemize}
\end{itemize}

\section{Features, Protocol, and the Leakage Audit}
\begin{itemize}
  \item \textbf{Sample.} 118 cells usable at $(W,H)=(8,4)$, status \texttt{ok}
        (54 phasic, 64 steady), from the 132-cell capture.
  \item \textbf{Features (10).} \texttt{apf\_mean}, \texttt{apf\_std},
        \texttt{cv\_workingset}, \texttt{f1\_phase}, \texttt{cepstral\_peak\_idx},
        \texttt{ceps\_peak\_snr\_db}, \texttt{stat\_pass\_frac},
        \texttt{n\_pairs}, \texttt{n\_windows}, \texttt{coverage\_ratio}.
  \item \textbf{Three are family-conditional} (the leakage vector):
    \begin{itemize}
      \item \texttt{cv\_workingset} is computed \emph{only} for steady cells
            (NaN for every phasic cell);
      \item \texttt{f1\_phase} is computed \emph{only} for phasic cells (NaN for
            every steady cell) -- complementary missingness that encodes the
            label;
      \item \texttt{coverage\_ratio} is a constant injective function of the
            family label: \emph{exactly} 0.1333 for every steady cell and 0.200
            for every phasic cell, zero within-family variance (verified in
            \texttt{plan03/sweep.csv}). It is the label re-encoded as a feature,
            so its 1.000 accuracy is definitional, not ``leakage to audit.''
    \end{itemize}
  \item \textbf{Protocol.} Median imputation $\to$ standardization $\to$
        classifier; GroupKFold with the replicate index as group
        (leave-one-replicate-out); models $k$-NN, RandomForest($n{=}300$), MLP.
  \item \textbf{Audit.} Re-run binary classification on feature subsets:
        full~10; family-agnostic (drop the three); each suspect feature alone;
        and a majority baseline. (Script: \texttt{leakage\_ablation.py};
        results: \texttt{leakage\_ablation.json}.)
\end{itemize}

\section{Results}
\begin{itemize}
  \item \textbf{The leakage is total.} Binary phasic-vs-steady accuracy
        (leave-one-replicate-out, RandomForest):
        \begin{center}
        \begin{tabular}{lc}
          \toprule
          Feature set & Accuracy \\
          \midrule
          Full (10 features)              & \textbf{1.000} \\
          \texttt{coverage\_ratio} alone  & \textbf{1.000} \\
          Family-agnostic (drop cv, f1, coverage) & 0.644 \\
          \texttt{ceps\_peak\_snr\_db} alone & 0.508 \\
          \texttt{apf\_mean}+\texttt{apf\_std} only & 0.576 \\
          Majority-class baseline         & 0.542 \\
          \bottomrule
        \end{tabular}
        \end{center}
  \item \textbf{Reading.} \texttt{coverage\_ratio} alone -- a feature the
        analyzer sets to 0.133 (steady) or 0.200 (phasic) by construction --
        reproduces the entire 100\%. The reported ``perfect family separation''
        is feature-construction leakage.
  \item \textbf{Residual (leakage-free) signal is weak -- and not shown to beat
        baseline.} With the three family-conditional features removed, accuracy
        is 0.644 against the 0.542 majority baseline (76 vs.\ 64 of 118 cells).
        No significance test is reported, and the gap does not survive the
        non-independence of the design: the four replicates per
        workload-duration are not independent, leaving $\approx 33$ independent
        groups, at which a one-sided proportion test gives $p \approx 0.12$
        (vs.\ $p \approx 0.013$ if the 118 cells were independent). The honest
        statement: a coarse scalar carries \emph{at most} a weak family hint,
        statistically indistinguishable from baseline at the replicate-group
        level. The paired test and grouped bootstrap are now in hand: an exact
        McNemar (agnostic vs.\ majority) gives $p=0.119$, and a 33-group cluster
        bootstrap puts agnostic accuracy at 0.645 with 95\% CI $[0.541, 0.746]$,
        whose lower bound coincides with the baseline -- so no ``genuine signal''
        claim is warranted (\texttt{robustness\_reanalysis.py}).
  \item \textbf{Cepstral SNR does not separate families.} Alone it scores 0.508
        (chance). Steady median SNR spans 2.65--7.99\,dB and overlaps the phasic
        range 4.78--6.35\,dB; the steady hash-table workload has the highest SNR
        of all. The first draft's ``phasic high / steady low SNR'' mechanism is
        retracted.
  \item \textbf{Confidence.} Even taken at face value, 118/118 has a
        Clopper--Pearson 95\% lower bound of 0.969 -- ``perfect'' is not
        defensible without an interval, and here it is not the right number
        anyway.
  \item \textbf{Instance (11-way).} RandomForest reaches 0.331 accuracy
        (macro-F1 0.328); $k$-NN and MLP 0.237. Best class
        \texttt{scanner\_metadata} F1 0.60; worst \texttt{workingset} 0.10,
        \texttt{mmap} 0.12. The block-diagonal-by-family structure of the 11-way
        matrix inherits the same leakage; the \emph{within}-family confusion is
        the genuine (and weak) part.
  \item \textbf{Re-test on a DOF-relieved capture (66 cells, 0 starved;}
        \textbf{delete-as-you-go pipeline, Paper~A/F).} The identical
        leakage-controlled protocol on a higher-fidelity capture, leakage-free
        (family-agnostic) RandomForest:
        \begin{center}
        \begin{tabular}{lcc}
          \toprule
          Protocol & v3 (starved) & relieved \\
          \midrule
          leave-one-replicate-out (LORO)        & 0.644 & 0.970 \\
          leave-one-workload-out (LOWO)         & 0.466 & 0.636 \\
          \quad + peak/variance features        & --    & 0.788 \\
          leave-one-\emph{family}-out (LOFO)    & --    & 0.30--0.64 \\
          anomaly ROC-AUC (novel family)        & --    & 0.93 \\
          majority baseline                     & 0.542 & 0.545 \\
          \bottomrule
        \end{tabular}
        \end{center}
  \item \textbf{Significance flips positive on relieved capture.} Exact McNemar
        (agnostic vs.\ majority) $p\approx0$; 33-group cluster bootstrap 95\% CI
        $[0.92, 1.0]$ -- against paperC's v3 $p=0.119$ with a CI lower bound at
        the baseline. The leakage-free family signal is real once DOF starvation
        is removed; relieving capture, not changing the model, moved it.
  \item \textbf{But it does not generalize across families.} Leave-one-family-out
        holding out all ransomware misses 3 of 4 subtypes (threat recall 0.25),
        and a novel subtype's family is recovered at chance -- a family is not a
        coherent cluster in APF. Peak/variance shape features lift the
        novel-\emph{workload} binary 0.636$\to$0.788 and the anomaly ROC-AUC
        0.74$\to$0.93, but do not confer cross-family generalization.
  \item \textbf{Normal-behavior baseline (benign cells, offline).} Treating normal
        behavior as a deliverable rather than a backdrop: a per-behavior
        fingerprint (median and p5--p95 of the leakage-free features) and a normal
        \emph{envelope} (per-feature benign p5--p95 band). An interpretable ``count
        features outside the envelope'' rule separates threat/benign at ROC-AUC
        0.95 (in-sample envelope, descriptive), complementing the held-out
        one-class detector (0.93). A \emph{masquerade map} -- the nearest benign
        behavior to each threat -- localizes the confusion: low-APF threats sit
        nearest benign \texttt{pagefault}, high-APF threats nearest \texttt{mmap},
        naming the specific normal behaviors each threat impersonates.
  \item TODO: figure -- accuracy vs.\ feature subset bar chart; table -- the
        per-family medians showing SNR overlap and the coverage-ratio constants.
\end{itemize}

\section{Limitations and the Security Question}
\begin{itemize}
  \item \textbf{Not a detector.} The ``phasic'' workloads are synthetic
        reversible-XOR sandboxes that emulate the memory-write \emph{rhythm} of
        encryption, not real malware; the ``steady'' class is memory
        micro-benchmarks, not benign applications. There is no benign-but-bursty
        control class, which is exactly where false positives would live. No
        security-detection claim is made or supported.
  \item \textbf{Required before any security framing} (future work): an explicit
        threat model and base rate; a real-malware corpus; a benign-phasic
        control; open-set / unknown-workload evaluation; ROC at a
        deployment-realistic base rate rather than balanced accuracy.
  \item \textbf{Statistical power -- now measured.} The leakage-free gap does not
        survive replicate non-independence. An exact McNemar test (family-agnostic
        vs.\ majority rule) gives $p=0.119$; a cluster bootstrap over the 33
        independent (workload, duration) groups puts agnostic accuracy at 0.645
        with 95\% CI $[0.541, 0.746]$, whose lower bound sits at the 0.542
        baseline. Under leave-one-\emph{workload}-out CV -- the harder test of
        generalizing to an unseen workload -- agnostic accuracy falls to 0.466,
        \emph{below} the majority baseline, while the leaky full set stays at
        1.000. This confirms both that \texttt{coverage\_ratio} is a per-family
        constant invariant to which workload is held out, and that the modest
        leave-one-replicate-out signal is partly workload memorization
        (\texttt{robustness\_reanalysis.py}).
  \item \textbf{Cross-family non-generalization persists on relieved capture.}
        The DOF-relieved re-test confirms the v3 leave-one-workload-out warning is
        structural, not a starvation artifact: leave-one-\emph{family}-out (hold
        out an entire family) misses 3 of 4 ransomware subtypes (recall 0.25), and
        a held-out subtype's family is recovered at chance. The
        leave-one-replicate-out gains are partly same-family memorization; APF does
        not cluster by family.
  \item \textbf{Anomaly detection is the regime that survives novelty.} A
        one-class detector trained only on benign flags a held-out threat family
        at ROC-AUC 0.93, because it models ``normal'' and never needs to have seen
        the threat -- the practical split is: known families $\to$ supervised
        classification; novel/unseen family $\to$ anomaly detection.
  \item \textbf{Imputation is a second family-conditional channel.} Because
        \texttt{cv\_workingset} and \texttt{f1\_phase} are missing by family
        (complementary NaN), median imputation injects a family-correlated
        constant and the missingness pattern itself standardizes differently per
        family. The family-agnostic run must \emph{drop} these columns entirely
        rather than impute-then-use; report a missingness-indicator ablation to
        confirm the channel is closed.
  \item \textbf{Scope.} Instance identity is not recoverable from APF; a per-page
        sub-signature (the 2-D pipeline of Paper A) would be required.
\end{itemize}

\section{Conclusion}
\begin{itemize}
  \item A single-scalar liveness trajectory carries \emph{at most} a weak family
        hint (leakage-free 0.644 vs.\ a 0.542 baseline, a gap not significant at
        the replicate-group level) and no usable instance signal; the apparent
        perfect separation is an artifact of family-conditional feature
        computation.
  \item Capture fidelity matters: on a DOF-relieved re-capture the same
        leakage-free protocol becomes significant (leave-one-workload-out 0.636
        vs.\ 0.545, McNemar $p\approx0$), but the signal still does not generalize
        across families (leave-one-family-out misses 3/4 of a held-out family),
        whereas a benign-only anomaly detector flags a novel family at ROC-AUC
        0.93. The family hint is real but coarse, and detection of an unseen
        family belongs to the anomaly regime, not the supervised one.
  \item The transferable lesson is methodological: when a pipeline derives
        per-class ``defining metrics,'' those metrics must be kept out of any
        classifier that is meant to demonstrate discrimination -- otherwise the
        model simply rediscovers the labelling rule.
\end{itemize}

\bibliographystyle{IEEEtran}
\bibliography{../refs}
\end{document}
